\documentclass[10pt, a4paper, twocolumn]{article}

% --- Pacotes ---
\usepackage[utf8]{inputenc}
\usepackage[T1]{fontenc}
\usepackage[brazil]{babel}
\usepackage[top=2cm, bottom=2cm, left=1.75cm, right=1.75cm]{geometry}
\usepackage{setspace}
\usepackage{indentfirst}
\usepackage{hyperref}
\usepackage{url}
\usepackage{enumitem}
\usepackage{booktabs}
\usepackage{graphicx}
\usepackage{float}
\usepackage{abstract}
\usepackage{amsmath}
\usepackage{multicol}
\usepackage{titlesec}

% --- Configurações de Espaçamento de Títulos (Otimização de espaço) ---
\titleformat{\section}{\normalfont\normalsize\bfseries}{\thesection}{1em}{}
\titleformat{\subsection}{\normalfont\normalsize\bfseries}{\thesubsection}{1em}{}
\titleformat{\subsubsection}{\normalfont\small\bfseries}{\thesubsubsection}{1em}{}
\titlespacing*{\section}{0pt}{1.2ex plus 0.5ex minus .2ex}{0.6ex plus .2ex}
\titlespacing*{\subsection}{0pt}{1.0ex plus 0.4ex minus .2ex}{0.5ex plus .1ex}
\titlespacing*{\subsubsection}{0pt}{0.9ex plus 0.3ex minus .2ex}{0.4ex plus .1ex}

% --- Configurações gerais ---
\setlength{\parindent}{0.4cm}
\singlespacing
\setlength{\columnsep}{0.5cm}

% Configuração para forçar quebra de linha em URLs longas nas referências
\makeatletter
\g@addto@macro{\UrlBreaks}{\do\/\do\-\do\.\do\=}
\makeatother

\hypersetup{
  colorlinks=true,
  linkcolor=black,
  urlcolor=blue,
  citecolor=black
}

% Estilização do Abstract Compacto
\renewenvironment{abstract}{
  \small
  \begin{center}
    \bfseries \abstractname\vspace{-0.6em}
  \end{center}
  \list{}{
    \setlength{\leftmargin}{0cm}
    \setlength{\rightmargin}{0cm}
  }%
  \item\relax
}
{\endlist}

% --- Título e autores ---
\title{
  \normalsize\textbf{Desenvolvimento e Análise de uma Ferramenta de Inteligência Artificial para Suporte Pedagógico Personalizado}
}

\author{
  \small \textbf{Raphael Victor de Araujo Alencar}, \textbf{Vinicius Haas Kuhn} \\[2pt]
  \small {Orientador: Prof. Me. Gustavo da Silva Andrade} \\[4pt]
  \small Universidade Católica Dom Bosco, Campo Grande, Brasil
}
\date{}

% ================================================================
\begin{document}

% Cabeçalho em coluna única integrada
\twocolumn[
  \begin{center}
    \maketitle
    \vspace{-2.2em}
    \begin{abstract}
Este trabalho apresenta o desenvolvimento, a especificação arquitetural e o plano de validação experimental de uma ferramenta de Inteligência Artificial (IA) baseada em aprendizagem adaptativa e suporte pedagógico personalizado para o ensino médio. O sistema visa mitigar a defasagem escolar por meio de auxílio mediado (\textit{scaffolding}), reduzindo o descarregamento cognitivo dos estudantes. A infraestrutura é organizada em três camadas independentes --- \textit{frontend} em React.js, \textit{backend} em FastAPI e um Orquestrador de IA acoplado à API do modelo Google Gemini Flash --- e foi materializada em uma implementação de referência funcional, com motor de adaptabilidade por janela móvel de cinco exercícios, autoria automática dos itens pelo próprio modelo de linguagem com controle explícito de não repetição e banco de reserva estático, dupla camada anti-vazamento de gabarito no modo socrático, defesa contra injeção de \textit{prompt}, limitação de taxa de chamadas e telemetria por tentativa. São descritas as alterações aplicadas à aplicação em relação à especificação inicial, entre elas a correção do limiar numérico da janela móvel e sua ampliação de três para cinco exercícios, a substituição do banco estático de questões pela geração sob demanda com \textit{fallback}, a retenção do gabarito exclusivamente no servidor e a substituição da garantia rígida de latência por degradação graciosa. Por fim, detalha-se o desenho da validação empírica por experimento controlado (teste A/B) com amostra projetada de 60 estudantes, cuja eficácia do suporte socrático face ao modo diretivo será quantificada pelo Ganho Normatizado de Hake, e discute-se a conclusão esperada do estudo, incluindo os resultados que confirmariam ou refutariam a hipótese central.

    \medskip
    \noindent\textbf{Palavras-chave:} Inteligência Artificial; Tutoria Inteligente; Aprendizagem Adaptativa; Defasagem Escolar; \textit{Scaffolding} Pedagógico; Modelos de Linguagem.
    \end{abstract}
    \vspace{1.2em}
  \end{center}
]

% ================================================================
\section{Introdução}

A pandemia de Covid-19 aprofundou uma defasagem educacional já estrutural no
Brasil. Segundo o relatório setorial documentado por Barbosa~\cite{cetic2025},
parcelas significativas de estudantes da educação básica acumularam lacunas de
aprendizagem que a retomada presencial, por si só, não foi suficiente para
reverter. Turmas numerosas, carga horária restrita e a heterogeneidade do ponto
de partida de cada aluno tornam inviável, para o professor humano, oferecer
acompanhamento verdadeiramente individualizado em larga escala. Nesse contexto,
cresce a pressão por soluções tecnológicas capazes de complementar o trabalho
docente sem substituí-lo.

A Inteligência Artificial (IA) tem emergido como uma dessas alternativas.
Sistemas de tutoria inteligente (\textit{Intelligent Tutoring Systems} --- ITS)
baseados em aprendizagem adaptativa já demonstram capacidade de mapear lacunas
específicas de cada estudante e propor trilhas de recuperação personalizadas,
conforme apontado por Vorobyeva et al.~\cite{roe2025}. Iniciativas como a
Estratégia Nacional de Escolas Conectadas do Ministério da
Educação~\cite{mec2023} sinalizam, por parte do poder público, o reconhecimento
de que a infraestrutura digital e o uso pedagógico da tecnologia são prioridades
para reduzir desigualdades educacionais. A IA, portanto, já não é uma
possibilidade futura, é uma realidade em expansão nas salas de aula brasileiras.

Entretanto, a forma como essa tecnologia tem sido incorporada revela um problema
central: predomina o uso substitutivo, em que o estudante delega à máquina a
produção da resposta em vez de utilizá-la como andaime para construir o próprio
raciocínio. Kestin et al.~\cite{kestin2025} identificam esse fenômeno como
``descarregamento cognitivo'' e demonstram que ele pode ser mitigado quando a IA
é programada para oferecer pistas e questionamentos socráticos em vez de
soluções prontas. A distinção entre um modo \textit{diretivo}, em que a
ferramenta entrega a resposta, e um modo \textit{socrático}, em que ela orienta o
estudante a chegar à resposta, é, portanto, central para o debate sobre o papel
pedagógico da IA.

Essa distinção, contudo, não é apenas uma escolha de redação de \textit{prompt}.
Modelos de linguagem de grande porte são sistemas estocásticos: mesmo instruídos
a não revelar a resposta, podem fazê-lo por descuido, por reformulação indireta
ou por indução deliberada do próprio estudante --- um caso particular de injeção
de \textit{prompt}. Um experimento que compare o modo socrático ao diretivo sem
mecanismos que garantam a separação entre os dois tratamentos corre o risco de
medir uma diferença que, na prática, não existiu. Esse ponto, pouco explorado na
literatura pedagógica, motivou a parte mais substantiva das alterações de
engenharia descritas neste trabalho.

Diante deste cenário, apresenta-se aqui a modelação técnica, a implementação de
referência e o plano experimental de uma ferramenta de IA voltada ao suporte
pedagógico individualizado no ensino médio. A contribuição é dupla: técnica, ao
especificar e construir uma arquitetura de três camadas com salvaguardas
explícitas de integridade do tratamento experimental; e pedagógica, ao formular
um protocolo rigoroso para testar a hipótese de que o auxílio mediado socrático
gera maior retenção de conhecimento do que o fornecimento direto de respostas.

Este artigo está organizado da seguinte forma: a Seção~\ref{sec:objetivos}
apresenta os objetivos geral e específicos; a Seção~\ref{sec:correlatos} discute
os trabalhos correlatos nas dimensões de \textit{scaffolding} pedagógico,
personalização em larga escala, integridade de sistemas baseados em modelos de
linguagem e integração ética; a Seção~\ref{sec:metodologia} detalha a arquitetura
do sistema, o motor de adaptabilidade, a autoria automática das questões, o
Orquestrador de IA, as alterações aplicadas à aplicação e o protocolo
experimental; a
Seção~\ref{sec:resultados} relata os resultados preliminares de validação
técnica; a Seção~\ref{sec:conclusao} discute a conclusão esperada; e a
Seção~\ref{sec:cronograma} apresenta o cronograma de execução.

% ================================================================
\section{Objetivos}
\label{sec:objetivos}

\subsection{Geral}

Desenvolver e validar experimentalmente uma ferramenta de Inteligência
Artificial voltada ao suporte pedagógico personalizado, capaz de mitigar a
defasagem de aprendizagem no ensino médio por meio de auxílio mediado,
avaliando seu impacto na autonomia de raciocínio dos estudantes.

\subsection{Específicos}

Para atingir o objetivo geral, foram estabelecidos os seguintes objetivos
específicos:

\begin{itemize}[leftmargin=0.4cm, itemsep=2pt, parsep=0pt]
  \item Projetar e implementar uma plataforma web de tutoria adaptativa que,
  a partir do desempenho do estudante em uma janela móvel de 5 exercícios,
  alterne dinamicamente o nível de dificuldade (básico, intermediário ou
  avançado) ou acione blocos de reforço conceitual.
  \item Implementar o módulo de processamento de linguagem natural integrado via
  API ao modelo Google Gemini Flash, configurado com engenharia de
  \textit{prompt} e saída estruturada para duas funções distintas: a autoria sob
  demanda de questões inéditas, ajustadas ao nível e ao estado de reforço do
  aluno e explicitamente instruídas a não repetir itens já respondidos, e a
  geração, imediatamente após o erro, de feedback diretivo ou socrático conforme
  o grupo experimental do estudante.
  \item Garantir a integridade do tratamento experimental por meio de
  mecanismos técnicos que impeçam o vazamento do gabarito no modo socrático,
  incluindo retenção do gabarito no servidor, autorrelato estruturado do modelo
  e filtro heurístico de pós-processamento.
  \item Instrumentar a aplicação com telemetria por tentativa (modo, latência,
  vazamento detectado, uso de \textit{fallback}) e com um painel de
  acompanhamento docente, permitindo tanto a análise estatística do experimento
  quanto a supervisão pedagógica humana.
  \item Comparar estatisticamente os modos de uso da ferramenta (diretivo versus
  socrático) por meio de experimento controlado (teste A/B), mensurando o ganho
  de aprendizagem por exames de pré e pós-teste.
  \item Conduzir análise dos riscos éticos da ferramenta com foco em plágio,
  desinformação, proteção de dados de menores e equidade de acesso, validando
  estratégias de mitigação técnica para uso institucionalizado.
\end{itemize}

% ================================================================
\section{Trabalhos Correlatos}
\label{sec:correlatos}

\subsection{IA como Andaime Pedagógico (\textit{Scaffolding})}

Estudos recentes têm investigado a transição da IA de ferramenta de busca
passiva para ``parceiro de tutoria interativo''. Kestin et
al.~\cite{kestin2025} demonstram que o uso de modelos de linguagem estruturados
com técnicas de \textit{scaffolding} (andaime) impede o chamado
``descarregamento cognitivo'' --- quando o aluno terceiriza o raciocínio para a
máquina. Em vez de fornecer a solução, a IA é programada para oferecer pistas e
questionamentos socráticos que guiam o estudante na construção do próprio
conhecimento.

Nesse sentido, Ocen et al.~\cite{ocen2026} realizaram um estudo de caso em
ambientes de aprendizagem aberta, identificando que a IA, quando configurada
como suporte ao aprendizado autodirigido, atua como mediadora que compensa a
falta de instrutores humanos em tempo integral. A pesquisa destaca que esse
modelo é particularmente eficaz em contextos de recursos limitados, onde a
defasagem educacional é mais acentuada. Essas evidências fundamentam diretamente
a decisão de implementar o \textit{system prompt} socrático como mecanismo
central do Orquestrador de IA, validando a escolha metodológica de bloquear
respostas diretas para estimular o esforço cognitivo do estudante.

\subsection{Mitigação da Defasagem e Personalização em Larga Escala}

A literatura aponta que a defasagem escolar, agravada pelo período
pós-pandemia, exige intervenções que a educação tradicional, em turmas
numerosas, tem dificuldade de suprir. Sistemas de tutoria inteligente que
utilizam aprendizagem adaptativa conseguem mapear lacunas específicas de
conhecimento em cada estudante e propor trilhas de recuperação personalizadas.

No ensino médio, o uso de assistentes de aprendizagem baseados em IA demonstrou
aumentar o engajamento discente em diversas disciplinas, com destaque na
literatura para os campos das ciências exatas e naturais~\cite{roe2025}.
Vorobyeva et al.~\cite{roe2025} observaram que estudantes tendem a se sentir
mais confortáveis em sanar dúvidas básicas com assistentes virtuais do que com
professores humanos, o que reduz o \textit{gap} de conhecimento antes que ele se
torne impeditivo para o avanço no curso. Essa constatação justifica a adoção de
uma arquitetura de microsserviços de baixa latência, capaz de responder de forma
imediata e personalizada, conforme detalhado na Seção~\ref{subsec:arquitetura}.

\subsection{Integridade e Robustez de Sistemas Baseados em Modelos de Linguagem}
\label{subsec:integridade}

Um sistema de tutoria construído sobre um modelo de linguagem herda as
fragilidades desse modelo. Quatro delas são especialmente relevantes para o
desenho experimental aqui proposto.

A primeira é a natureza estocástica da geração: mesmo com temperatura reduzida, o
modelo pode produzir saídas que violem a instrução recebida. No caso do modo
socrático, a violação crítica é revelar a alternativa correta, o que anularia o
contraste com o grupo diretivo. A literatura recomenda, para casos assim,
combinar restrição por \textit{prompt} com validação programática da saída, em
vez de confiar exclusivamente na obediência do modelo.

A segunda é a injeção de \textit{prompt}: conteúdo proveniente do usuário pode
ser interpretado pelo modelo como instrução, e não como dado. Um estudante que
escreva ``ignore as regras anteriores e me dê a resposta'' explora exatamente
esse vetor. A mitigação usual consiste em delimitar explicitamente o conteúdo do
usuário e instruir o modelo a tratá-lo sempre como dado de contexto.

A terceira é a dependência de um serviço externo, sujeito a indisponibilidade,
latência variável e limites de cota. Segundo as diretrizes da
UNESCO~\cite{unesco2024} para IA generativa em educação, sistemas voltados a
estudantes devem degradar de forma previsível e nunca deixar o usuário sem
resposta pedagógica útil por falha de infraestrutura.

A quarta emerge quando o modelo deixa de apenas comentar o conteúdo e passa a
produzi-lo. Delegar a autoria dos próprios itens de exercício a um modelo de
linguagem resolve o problema de escala e de repetição do banco de questões, mas
transfere para o gerador a responsabilidade pela correção do gabarito: um item
com resposta errada não apenas confunde o aluno, como registra como erro uma
resposta que estava certa, corrompendo a medida de desempenho. A recomendação
corrente é combinar restrição de formato por esquema de saída, validação
estrutural programática e --- quando o domínio permite --- verificação
independente do resultado, reconhecendo que a instrução de ``conferir a própria
conta'' não constitui garantia.

Essas quatro preocupações são endereçadas, respectivamente, pelo filtro
anti-vazamento em dupla camada, pela delimitação do conteúdo do aluno, pelo
mecanismo de \textit{fallback} descritos na Seção~\ref{subsec:orquestrador} e
pela validação estrutural da questão gerada, com recurso ao banco de reserva,
descrita na Seção~\ref{subsec:autoria}.

\subsection{Integração Ética e o Futuro da Práxis Docente}

Diferentemente da abordagem de automação pura, a tendência atual na pesquisa
acadêmica defende o modelo ``IA gerenciada pelo professor''. Pesquisas recentes
compararam instruções lideradas por docentes com e sem auxílio de IA, concluindo
que o melhor desempenho ocorre quando o professor atua como arquiteto da
experiência, utilizando a IA para feedback imediato e análise de dados, enquanto
mantém a liderança no desenvolvimento do pensamento crítico e
ético~\cite{luckin2018}. A partir desse entendimento, a ferramenta incorpora um
painel administrativo que expõe as métricas agregadas do experimento e permite ao
docente reatribuir manualmente o modo de mediação de um aluno específico,
unindo a automação técnica à supervisão pedagógica humana.

% ================================================================
\section{Materiais e Métodos}
\label{sec:metodologia}

Esta seção descreve o público-alvo, a arquitetura implementada, o motor de
adaptabilidade, o Orquestrador de IA, as alterações aplicadas à aplicação em
relação à especificação inicial e o protocolo experimental desenhado para a
validação em campo.

\subsection{Delimitação do Público-Alvo e Critérios de Inclusão}
\label{subsec:publico_alvo}

A validação prática e a coleta de dados empíricos serão realizadas junto a
estudantes ingressantes do Ensino Médio, com foco em módulos de nivelamento.
Para a composição da amostra, serão adotados os seguintes critérios de inclusão:

\begin{itemize}[leftmargin=0.4cm, itemsep=2pt, parsep=0pt]
    \item Estar regularmente matriculado no 1º ano do Ensino Médio da
    instituição de ensino parceira;
    \item Obter autorização formal dos pais ou responsáveis legais mediante
    assinatura do Termo de Consentimento Livre e Esclarecido (TCLE);
    \item Manifestar concordância voluntária por meio da assinatura do Termo de
    Assentimento Livre e Esclarecido (TALE) para menores;
    \item Realizar todas as etapas obrigatórias do protocolo, que incluem o teste
    de diagnóstico inicial (pré-teste) e a avaliação somativa final (pós-teste).
\end{itemize}

Cabe registrar uma delimitação de escopo: embora a proposta original previsse
módulos de Matemática e de Língua Portuguesa, a implementação cobre até o
momento apenas conteúdos de Matemática. O domínio é delimitado por uma lista de
assuntos associada a cada nível, que restringe tanto a geração de questões pela
IA quanto o banco de reserva estático: aritmética básica, frações, porcentagem,
equação do 1º grau e regra de três simples no nível básico; equação do 2º grau,
sistemas lineares $2 \times 2$, função afim e progressão aritmética no
intermediário; função exponencial, logaritmos, progressão geométrica e derivadas
no avançado. O módulo de Língua Portuguesa permanece como trabalho futuro, e a
análise estatística será conduzida sobre o domínio efetivamente coberto.

\subsection{Arquitetura do Sistema}
\label{subsec:arquitetura}

A ferramenta está estruturada em três camadas independentes, projetadas para
isolar responsabilidades e permitir que falhas na dependência externa não
comprometam o restante do fluxo pedagógico:

\begin{enumerate}[leftmargin=0.4cm, itemsep=3pt]
    \item \textbf{Camada Cliente (\textit{Frontend}):} interface reativa
    desenvolvida em React.js com TypeScript, organizada como \textit{Single Page
    Application} (SPA), com as telas de autenticação, diagnóstico, exercício,
    painel do aluno e painel administrativo. A comunicação com o servidor ocorre
    por HTTP com token JWT.

    \item \textbf{Núcleo de Negócio (\textit{Backend}):} API assíncrona
    construída em Python com o \textit{framework} FastAPI, responsável pelas
    regras de adaptabilidade, controle de sessão via JWT, persistência
    relacional através de SQLAlchemy e validação das respostas. Toda a
    verificação de acerto ocorre no servidor, contra o gabarito armazenado no
    banco.

    \item \textbf{Orquestrador de IA:} módulo isolado que concentra as chamadas à
    API do Google Gemini Flash --- tanto a autoria das questões quanto a geração
    do feedback ---, a montagem dos \textit{system prompts}, a validação da saída
    e a política de \textit{fallback}. Nenhuma outra parte da aplicação conversa
    diretamente com o modelo.
\end{enumerate}

A Tabela~\ref{tab:stack} resume a pilha tecnológica e os parâmetros de
configuração adotados.

\begin{table}[H]
\centering
\small
\caption{Pilha tecnológica e parâmetros de configuração}
\label{tab:stack}
\vspace{0.13cm}
\resizebox{\columnwidth}{!}{
\begin{tabular}{ll}
\toprule
\textbf{Componente} & \textbf{Tecnologia / valor} \\ \midrule
Interface & React.js + TypeScript (SPA, Vite) \\
API & FastAPI (Python), SQLAlchemy \\
Persistência & SQLite (desenvolvimento) / PostgreSQL \\
Autenticação & JWT (HS256), expiração de 120 min \\
Hash de senha & bcrypt com \textit{salt} (\textit{passlib}) \\
Modelo de linguagem & Google Gemini Flash \\
Temperatura & $T = 0{,}1$ (feedback) / $0{,}2$ (autoria) \\
Saída do modelo & JSON validado por esquema \\
Limite de chamadas & 30 por usuário por minuto \\
Janela de adaptação & 5 exercícios, sem sobreposição \\
Níveis & básico, intermediário, avançado \\
Origem das questões & geradas pela IA sob demanda \\
Histórico anti-repetição & 20 enunciados mais recentes \\
Banco de reserva & 15 itens (3 diagnósticos, 4 de reforço) \\ \bottomrule
\end{tabular}%
}
\end{table}

\subsection{Motor de Adaptabilidade}
\label{subsec:adaptabilidade}

A adaptação da dificuldade é regida por uma janela móvel de $W = 5$ exercícios,
avaliada apenas quando o bloco fecha, sem sobreposição de margens. Seja $c$ o
número de acertos dentro da janela corrente e $L = \{$básico, intermediário,
avançado$\}$ a sequência ordenada de níveis, a transição de estado obedece à
regra da Equação~\ref{eq:janela}:

\begin{equation}
\text{ação} =
\begin{cases}
\text{promover nível,} & c = W \quad (100\%) \\
\text{manter nível,}   & \lfloor W/3 \rfloor < c < W \\
\text{acionar reforço,} & c \leq \lfloor W/3 \rfloor
\end{cases}
\label{eq:janela}
\end{equation}

Com $W = 5$, o limiar de reforço é $\lfloor 5/3 \rfloor = 1$ acerto: promove-se o
aluno com 5 acertos (100\%), aciona-se reforço com no máximo 1 acerto
($\leq 20\%$) e mantém-se o nível na faixa intermediária de 2 a 4 acertos (40\% a
80\%). O limiar é expresso em função de $W$ justamente para que a proporção de
desempenho que dispara o reforço --- em torno de um terço de aproveitamento ---
permaneça estável caso o tamanho da janela seja reajustado.

Ao fim de cada avaliação, os contadores da janela são zerados. Quando o estado de
reforço está ativo, a questão seguinte é gerada (ou, no \textit{fallback},
selecionada) no formato remedial para o nível corrente do aluno, retornando ao
fluxo regular no bloco seguinte. A Tabela~\ref{tab:janela} sintetiza a regra e a
compara com a formulação da especificação inicial, cuja correção é discutida na
Seção~\ref{subsec:alteracoes}.

\begin{table}[H]
\centering
\small
\caption{Regra da janela móvel de 5 exercícios}
\label{tab:janela}
\vspace{0.13cm}
\resizebox{\columnwidth}{!}{
\begin{tabular}{lll}
\toprule
\textbf{Acertos} & \textbf{Aproveitamento} & \textbf{Ação do sistema} \\ \midrule
5 de 5 & 100\% & Promove ao próximo nível \\
4 de 5 & 80\%  & Mantém nível (banda de transição) \\
3 de 5 & 60\%  & Mantém nível (banda de transição) \\
2 de 5 & 40\%  & Mantém nível (banda de transição) \\
1 de 5 & 20\%  & Aciona bloco de reforço conceitual \\
0 de 5 & 0\%   & Aciona bloco de reforço conceitual \\ \bottomrule
\end{tabular}%
}
\end{table}

O nível inicial é atribuído por um teste de nivelamento de 3 problemas básicos.
O mapeamento é deliberadamente conservador: o aluno só é alocado diretamente no
nível intermediário com 100\% de acerto; qualquer erro o mantém no nível básico.
A pequena dimensão desse teste ($n = 3$) é uma limitação metodológica
reconhecida --- é uma amostra fraca para estimar proficiência --- mitigada pelo
fato de que a janela móvel reavalia o aluno continuamente a cada cinco
exercícios, de modo que uma alocação inicial imprecisa se corrige em poucos
blocos.

\subsection{Autoria de Questões pelo Modelo de Linguagem}
\label{subsec:autoria}

Na especificação inicial, os exercícios eram retirados de um banco estático
previamente redigido. Essa solução mostrou-se incompatível com a própria lógica
adaptativa: com poucos itens por nível, um aluno que permanecesse vários blocos
no mesmo estágio reencontrava questões já respondidas, e o reencontro contamina
a medida --- o acerto passa a refletir memorização do item, e não domínio do
conceito. A implementação atual, portanto, passou a gerar cada exercício sob
demanda por meio do próprio modelo de linguagem, mantendo o banco estático
apenas como reserva.

A cada requisição de novo exercício, o Orquestrador é acionado com três
parâmetros de estado: o nível corrente do aluno, o indicador de reforço e a
lista dos enunciados que ele já respondeu. Um \textit{system prompt} específico
de autoria instrui o modelo a escrever uma única questão inédita de múltipla
escolha, no estilo de vestibular, com exatamente quatro alternativas
identificadas por \textit{a}, \textit{b}, \textit{c} e \textit{d}, das quais
apenas uma correta e as demais correspondentes a erros plausíveis de raciocínio;
a conferir o próprio cálculo antes de responder; e a redigir, no mesmo retorno, a
explicação completa da resolução, armazenada para uso posterior caso o aluno
erre. O assunto é sorteado da lista associada ao nível
(Seção~\ref{subsec:publico_alvo}), o que impede que o gerador concentre os itens
em um único tópico. Quando o estado de reforço está ativo, uma instrução
adicional pede um item mais simples e guiado, centrado em um único conceito de
base em vez de um problema aplicado.

O controle de repetição é explícito: os 20 enunciados mais recentes do aluno são
enviados no \textit{prompt} sob instrução de não os repetir nem produzir
variação próxima, e o enunciado devolvido é ainda comparado programaticamente a
esse conjunto, sendo descartado em caso de coincidência. A saída é forçada em
JSON por esquema de resposta e validada estruturalmente após o \textit{parse}:
verifica-se que há exatamente quatro alternativas, que seus identificadores são
únicos e correspondem ao conjunto esperado, e que o gabarito declarado aponta
para uma delas. A temperatura é mantida em $0{,}2$, valor baixo o bastante para
favorecer consistência aritmética sem colapsar a variedade dos enunciados.

Sobre esse caminho opera uma política de \textit{fallback} análoga à do
feedback: se a chamada falhar, se a saída violar a validação estrutural ou se o
modelo insistir em repetir uma questão já vista, o sistema recorre ao banco de
reserva estático, preferindo itens ainda não respondidos pelo aluno. A decisão
de projeto é deliberada --- a disponibilidade da atividade pedagógica importa
mais do que o ineditismo do item ---, e o episódio é registrado. Cada questão
gerada é persistida com marcação de origem e vinculada ao aluno para o qual foi
criada, o que preserva a rastreabilidade necessária à auditoria posterior das
tentativas.

\subsection{Orquestrador de IA e Salvaguardas}
\label{subsec:orquestrador}

Além da autoria dos itens descrita na Seção~\ref{subsec:autoria}, o Orquestrador
é o componente responsável por transformar um erro do aluno em feedback
pedagógico. Nessa segunda função, ele opera em dois modos, selecionados pelo
grupo experimental ao qual o estudante foi sorteado:

\begin{itemize}[leftmargin=0.4cm, itemsep=3pt]
    \item \textbf{Modo Diretivo (grupo controle):} o \textit{system prompt}
    instrui o modelo a explicar a resolução completa e correta do problema,
    passo a passo, incluindo a resposta final.
    \item \textbf{Modo Socrático (grupo experimental):} o \textit{system prompt}
    proíbe explicitamente revelar a alternativa correta, mesmo de forma parcial
    ou indireta. O modelo deve identificar o provável erro de raciocínio por trás
    da alternativa escolhida e devolver de uma a duas perguntas orientadoras.
\end{itemize}

Sobre esses dois modos operam quatro salvaguardas encadeadas, motivadas pelas
fragilidades discutidas na Seção~\ref{subsec:integridade}:

\begin{description}[style=nextline, leftmargin=0cm, itemsep=5pt]
    \item[\textbf{1. Retenção do gabarito no servidor}]
    O identificador da alternativa correta nunca é serializado para o cliente
    antes da submissão. Após a submissão, ele é devolvido apenas aos alunos do
    grupo diretivo; no grupo socrático, o campo permanece nulo. Sem essa medida,
    o vazamento ocorreria na própria resposta HTTP, tornando irrelevante
    qualquer cuidado com o texto gerado pelo modelo.

    \item[\textbf{2. Saída estruturada com autorrelato}]
    A chamada ao modelo é configurada com esquema de resposta em JSON contendo
    dois campos: o conteúdo do feedback e um booleano em que o próprio modelo
    declara se acabou revelando a alternativa correta. Isso converte uma
    instrução em linguagem natural em um sinal verificável programaticamente.

    \item[\textbf{3. Filtro heurístico de pós-processamento}]
    Independentemente do autorrelato, o texto gerado é normalizado (remoção de
    acentuação, caixa baixa e colapso de pontuação) e comparado ao texto do
    gabarito, além de ser inspecionado por expressão regular à procura de
    menções do tipo ``alternativa X'' ou ``letra X'' quando $X$ é a resposta
    correta. Trata-se de defesa em profundidade: o autorrelato falha quando o
    modelo não percebe o próprio vazamento; a heurística falha diante de
    paráfrases; a probabilidade de ambos falharem simultaneamente é
    substancialmente menor.

    \item[\textbf{4. \textit{Fallback} pedagógico}]
    Se qualquer uma das duas camadas anteriores acusar vazamento, a saída do
    modelo é descartada e substituída por um andaime genérico pré-redigido, que
    devolve ao aluno perguntas de reexame do próprio raciocínio. O mesmo
    mecanismo é acionado quando a chamada à API falha ou excede o tempo
    tolerável. Em ambos os casos, o evento é registrado na telemetria.
\end{description}

Adicionalmente, o conteúdo proveniente do aluno é encapsulado entre delimitadores
explícitos e rotulado, no próprio \textit{prompt}, como dado de contexto e nunca
como instrução --- incluindo a instrução expressa de ignorar qualquer comando
nele contido, mesmo que o remetente alegue ser professor ou administrador. Por
fim, um limitador de taxa restringe as chamadas ao modelo a 30 por usuário por
minuto, protegendo o custo da API e contendo tentativas de exaustão do serviço.
O teto foi elevado em relação ao valor originalmente adotado porque a geração de
questões passou a consumir a mesma cota do feedback: cada exercício exibido é,
agora, uma chamada adicional ao modelo, e um limite baixo demais interromperia o
aluno apenas por avançar entre questões.

\subsection{Fluxo Operacional da Plataforma}
\label{subsec:fluxo}

O funcionamento do sistema ocorre em seis fases sequenciais:

\begin{description}[style=nextline, leftmargin=0cm, itemsep=5pt]
    \item[\textbf{Passo 1: Cadastro, sorteio e diagnóstico}]
    No cadastro, o estudante é sorteado de forma cega para um dos dois grupos
    experimentais, em proporção 1:1. Em seguida responde ao teste de nivelamento
    de 3 problemas básicos, do qual resulta sua assinatura inicial de
    proficiência.

    \item[\textbf{Passo 2: Geração adaptativa da questão}]
    O motor de adaptabilidade consulta o estado atual do aluno --- nível,
    indicador de reforço e histórico de enunciados já respondidos --- e solicita
    ao Orquestrador uma questão inédita para esse estado. Havendo falha,
    saída inválida ou repetição, recorre-se ao banco de reserva estático.

    \item[\textbf{Passo 3: Submissão e validação}]
    O aluno submete a alternativa escolhida. O servidor compara a submissão com
    o gabarito armazenado no banco e registra a tentativa.

    \item[\textbf{Passo 4: Avaliação da janela móvel}]
    A cada bloco fechado de 5 exercícios, aplica-se a regra da
    Equação~\ref{eq:janela}, promovendo o aluno, mantendo-o no nível ou
    acionando reforço conceitual.

    \item[\textbf{Passo 5: Geração de feedback}]
    Havendo erro, o Orquestrador de IA é acionado no modo correspondente ao grupo
    do aluno, com todas as salvaguardas da Seção~\ref{subsec:orquestrador}
    ativas.

    \item[\textbf{Passo 6: Persistência de telemetria}]
    Cada tentativa é registrada com nível vigente, alternativa escolhida e
    resultado; cada chamada de \textit{feedback} ao modelo gera um registro
    adicional com modo, latência medida, ocorrência de vazamento detectado, uso
    de \textit{fallback} e um trecho da saída bruta para auditoria posterior. As
    questões geradas ficam persistidas com marcação de origem e vínculo ao aluno,
    permitindo reconstruir a posteriori o item exato que cada tentativa
    respondeu.
\end{description}

\subsection{Alterações Realizadas na Aplicação}
\label{subsec:alteracoes}

A passagem da especificação para a implementação funcional expôs
inconsistências e lacunas que exigiram alterações de projeto. Esta seção as
documenta, por transparência metodológica e porque várias delas afetam
diretamente a validade do experimento. A Tabela~\ref{tab:alteracoes} sintetiza
as mudanças; o texto a seguir detalha as de maior impacto.

\begin{table}[H]
\centering
\small
\caption{Alterações aplicadas em relação à especificação inicial}
\label{tab:alteracoes}
\vspace{0.13cm}
\resizebox{\columnwidth}{!}{
\begin{tabular}{ll}
\toprule
\textbf{Ponto da especificação} & \textbf{Alteração aplicada} \\ \midrule
Limiar ``$\leq$66\% = 2 erros em 3'' & Regra de 3 faixas em função de $W$ \\
Janela móvel de 3 exercícios & Janela ampliada para $W = 5$ \\
Banco estático de questões & Autoria pela IA + banco de reserva \\
Repetição de itens não tratada & Histórico de 20 enunciados no \textit{prompt} \\
Transição só ``básico $\rightarrow$ interm.'' & Ordem genérica de níveis, promoção em cadeia \\
Gabarito no cliente & Nunca serializado antes da submissão \\
Filtro anti-vazamento vago & Dupla camada: autorrelato + heurística \\
Injeção de \textit{prompt} não tratada & Conteúdo do aluno delimitado como dado \\
Sem telemetria & Registro por tentativa e por chamada ao modelo \\
Sem controle de custo & Limite de 30 chamadas/usuário/minuto \\
``Latência garantida $<$ 3 s'' & Degradação graciosa com \textit{fallback} \\
Anonimização por SHA-256 & bcrypt com \textit{salt} para senhas \\
Modelo Gemini 1.5 Flash & Gemini Flash em versão corrente, configurável \\
Professor sem controle & Reatribuição manual de grupo pelo painel \\ \bottomrule
\end{tabular}%
}
\end{table}

\subsubsection{Correção do limiar e ampliação da janela móvel}

A especificação inicial definia o gatilho de reforço como ``desempenho igual ou
inferior a 66\% (2 erros em 3 questões)''. A equivalência está incorreta: 2 erros
em 3 correspondem a 33\% de aproveitamento, não a 66\%; 66\% é o resultado de 2
\textit{acertos} em 3. Implementar a regra literalmente faria com que alunos com
desempenho razoável fossem enviados a reforço conceitual, distorcendo a trilha de
aprendizagem de ambos os grupos. A formulação foi substituída pela regra de três
faixas da Equação~\ref{eq:janela}, que além de corrigir o cálculo resolve o caso
intermediário que a redação original deixava indefinido.

Em um segundo momento, o tamanho da janela foi ampliado de 3 para 5 exercícios.
A motivação é estatística: com $W = 3$, a decisão de promover ou rebaixar o aluno
repousava sobre três observações, e um único acerto ou erro por acaso ---
chute em questão de múltipla escolha com quatro alternativas, distração pontual
--- deslocava a trilha inteira. Com $W = 5$, a mesma decisão passa a se apoiar em
cinco observações, reduzindo a variância da estimativa de proficiência e tornando
as transições de nível menos oscilantes, ao custo de um tempo de resposta um
pouco maior do sistema a mudanças reais de desempenho. O limiar de reforço foi
reescrito como $\lfloor W/3 \rfloor$ para preservar a proporção original ---
aproximadamente um terço de aproveitamento --- em vez de fixar um número
absoluto de acertos, de modo que a regra permaneça coerente sob qualquer novo
ajuste de $W$. Nota-se que a ampliação também alarga a banda de manutenção de
nível, que passa a cobrir de 40\% a 80\% de aproveitamento: a promoção continua
exigindo desempenho perfeito no bloco, critério deliberadamente conservador,
coerente com o adotado no teste de nivelamento.

\subsubsection{Substituição do banco estático pela geração sob demanda}

A especificação previa um banco de questões previamente redigido, do qual o
motor sortearia o próximo item. Como discutido na Seção~\ref{subsec:autoria},
essa estratégia produz reencontro de itens já respondidos assim que o aluno
permanece alguns blocos no mesmo nível --- efeito que a ampliação da janela para
5 exercícios, por consumir mais itens por bloco, só faria acentuar. Como o acerto
em uma questão já vista mede memorização e não domínio conceitual, o problema não
é apenas de experiência de uso: ele contamina a variável de acurácia durante a
intervenção, que é uma das métricas do experimento.

A implementação passou, então, a solicitar ao modelo de linguagem a autoria de
cada exercício, informando o nível, o estado de reforço e o histórico recente de
enunciados a evitar. O banco previamente redigido foi rebaixado à condição de
reserva, acionado apenas quando a geração falha, devolve estrutura inválida ou
reincide em um enunciado já visto. Trata-se da alteração de maior alcance
funcional desta revisão, e ela desloca uma limitação reconhecida nos resultados
preliminares --- a insuficiência do banco para uma coleta real --- para um novo
ponto de atenção, a confiabilidade matemática do item gerado, discutido na
Seção~\ref{sec:resultados}.

\subsubsection{Generalização das transições de nível}

A especificação descrevia apenas a transição ``básico $\rightarrow$
intermediário'', sem definir o comportamento nos demais níveis. A implementação
adota uma sequência ordenada de níveis percorrida por índice, de modo que a
promoção funcione de forma idêntica em toda a escala e sature no nível avançado.

\subsubsection{Retenção do gabarito e integridade do experimento}

A especificação não tratava do momento em que o gabarito chega ao cliente. Como
discutido na Seção~\ref{subsec:orquestrador}, entregá-lo junto com o enunciado
--- prática comum em aplicações que validam respostas no navegador ---
invalidaria o tratamento socrático por completo. A validação foi integralmente
deslocada para o servidor e o gabarito passou a ser devolvido apenas ao grupo
diretivo, após a submissão.

\subsubsection{Substituição da garantia de latência}

A especificação afirmava que o sistema garantiria resposta em menos de 2 a 3
segundos. Trata-se de uma promessa que nenhuma aplicação pode fazer sobre uma
dependência externa de terceiros. A garantia rígida foi substituída por um
requisito de degradação graciosa: a chamada ao modelo é isolada e, em caso de
falha ou lentidão, o aluno recebe um andaime pedagógico genérico em vez de uma
tela travada. A latência real de cada chamada passou a ser medida e registrada,
convertendo-se de promessa em variável observável do experimento.

\subsubsection{Correção da estratégia criptográfica}

A especificação previa ``anonimização das credenciais via criptografia
SHA-256''. SHA-256 é uma função de resumo de propósito geral e rápida, imprópria
para armazenamento de senhas justamente por sua velocidade. A implementação
utiliza bcrypt com \textit{salt}, algoritmo de derivação deliberadamente custoso,
que é a prática recomendada para essa finalidade. A pseudonimização dos registros
de interação para fins de análise permanece conforme previsto.

\subsubsection{Instrumentação e supervisão docente}

Foram acrescentados dois elementos ausentes da especificação. O primeiro é a
telemetria por tentativa, sem a qual não haveria como calcular a taxa de
vazamento nem a latência por grupo. O segundo é o painel administrativo, que
expõe as métricas agregadas por grupo e permite ao docente reatribuir
manualmente o modo de mediação de um aluno --- concretizando o modelo de ``IA
gerenciada pelo professor'' discutido na revisão da literatura.

\subsection{Protocolo Experimental}
\label{subsec:experimento}

\subsubsection{Amostra, randomização e aspectos éticos}

O experimento será conduzido com amostra intencional de $N = 60$ estudantes
regularmente matriculados no 1º ano do Ensino Médio. No cadastro, cada usuário é
alocado de forma cega e automatizada, em proporção 1:1, por rotina de
randomização computacional:

\begin{itemize}[leftmargin=0.4cm, itemsep=2pt]
    \item \textbf{Grupo A (controle --- modo diretivo):} $n_A \approx 30$
    estudantes recebem resoluções diretas e imediatas geradas pela IA após o
    erro.
    \item \textbf{Grupo B (experimental --- modo socrático):} $n_B \approx 30$
    estudantes recebem pistas conceituais e questionamentos reflexivos.
\end{itemize}

Como o sorteio é independente por cadastro, os tamanhos finais dos grupos podem
divergir levemente de 30; o painel administrativo permite ao pesquisador
reequilibrar a alocação, se necessário, antes do início da coleta.

Em cumprimento à Lei Geral de Proteção de Dados (LGPD) no tratamento de dados de
menores, o sistema adota o princípio da minimização, coletando apenas dados
estritamente necessários ao escopo pedagógico: respostas, tempos de interação e
metadados das chamadas ao modelo. As senhas são armazenadas sob bcrypt com
\textit{salt} e os registros operacionais passam por pseudonimização na base. A
política de retenção é de 6 meses --- tempo de execução e defesa do projeto ---
seguida de descarte definitivo. O acesso à base bruta é restrito aos
pesquisadores responsáveis e protegido por perfil administrativo distinto do
perfil de aluno. A execução está condicionada à coleta prévia do TCLE junto aos
responsáveis e do TALE junto aos alunos.

\subsubsection{Métricas coletadas}

A Tabela~\ref{tab:metricas} lista as variáveis registradas e sua função na
análise.

\begin{table}[H]
\centering
\small
\caption{Variáveis coletadas e finalidade}
\label{tab:metricas}
\vspace{0.13cm}
\resizebox{\columnwidth}{!}{
\begin{tabular}{ll}
\toprule
\textbf{Variável} & \textbf{Finalidade} \\ \midrule
Ganho de Hake ($g$) & Variável dependente principal \\
Acurácia por grupo & Desempenho durante a intervenção \\
Distribuição de níveis & Progressão na trilha adaptativa \\
Latência do modelo (ms) & Viabilidade operacional \\
Taxa de vazamento detectado & Integridade do tratamento socrático \\
Uso de \textit{fallback} & Robustez frente a falhas \\
Tentativas até acerto & Indicador de esforço cognitivo \\ \bottomrule
\end{tabular}%
}
\end{table}

A taxa de vazamento detectado merece destaque: ela não é apenas um indicador de
qualidade de software, mas uma \textit{medida de validade interna do
experimento}. Uma taxa elevada indicaria que parte dos alunos do grupo socrático
recebeu, na prática, tratamento diretivo, e exigiria análise de sensibilidade
excluindo essas interações.

\subsubsection{Formulação matemática do ganho de aprendizagem}

Para testar a hipótese de mitigação do descarregamento cognitivo, aplicar-se-á a
métrica de Ganho Normatizado de Hake ($g$), calculada individualmente a partir do
desempenho em pré-teste (diagnóstico) e pós-teste (somativo, sem auxílio da IA),
conforme a Equação~\ref{eq:ganho_hake}:

\begin{equation}
    g = \frac{\text{\% Pós} - \text{\% Pré}}{100\% - \text{\% Pré}}
    \label{eq:ganho_hake}
\end{equation}

Onde $\text{\% Pré}$ corresponde ao percentual de acertos na avaliação de base e
$\text{\% Pós}$ ao aproveitamento no teste de retenção final. Após verificação de
normalidade pelo teste de Shapiro-Wilk, os valores médios de $g$ dos dois grupos
serão comparados pelo teste $t$ de Student para amostras independentes, com nível
de significância bicaudal $\alpha = 0{,}05$. Caso a normalidade seja rejeitada,
adotar-se-á o teste não paramétrico de Mann-Whitney como alternativa. Reporta-se
também o tamanho de efeito ($d$ de Cohen), dado que, com $N = 60$, um efeito
pequeno pode não atingir significância estatística sem que isso implique
ausência de efeito real.

% ================================================================
\section{Resultados Preliminares}
\label{sec:resultados}

Até o momento, os resultados obtidos são de natureza técnica: nenhuma coleta com
estudantes foi realizada, e portanto nenhum dado pedagógico é reportado aqui.

A implementação de referência das três camadas está funcional. O fluxo completo
--- cadastro com sorteio de grupo, teste de nivelamento, geração dos exercícios
pelo modelo, ciclo de resposta, promoção de nível, acionamento de reforço e
geração de feedback nos dois modos --- foi exercitado ponta a ponta contra o
servidor em execução. A compilação do
\textit{frontend} (verificação de tipos e empacotamento de produção) conclui sem
erros.

Três comportamentos das salvaguardas foram verificados manualmente. Primeiro, a
ausência do gabarito nas respostas HTTP anteriores à submissão, confirmando o
isolamento descrito na Seção~\ref{subsec:orquestrador}. Segundo, a degradação
graciosa: sem chave de API configurada, o sistema não interrompe o fluxo, entrega
o andaime genérico e registra o evento como uso de \textit{fallback}. Terceiro, a
substituição da saída do modelo quando o filtro heurístico acusa menção direta à
alternativa correta.

O caminho de autoria de questões foi igualmente exercitado. Verificou-se que os
itens gerados chegam no formato esperado --- quatro alternativas com
identificadores únicos e gabarito apontando para uma delas ---, que enunciados
estruturalmente inválidos são rejeitados antes de chegarem ao aluno, e que a
ausência de chave de API leva o sistema a servir o banco de reserva em vez de
interromper a atividade. O banco de reserva contém 15 itens, sendo 3 de
diagnóstico e 4 de reforço conceitual, distribuídos pelos três níveis. Esse
volume, que era insuficiente enquanto constituía a fonte primária de exercícios,
é adequado à função atual: cobrir os intervalos em que a geração falha, até que
a chamada seguinte seja bem-sucedida.

Três limitações técnicas permanecem em aberto e estão registradas por
transparência. A primeira, e a mais relevante para a validade da coleta, é que a
validação da questão gerada é estrutural, não matemática: o sistema garante o
formato do item, mas não verifica de forma independente se o gabarito declarado
é de fato a resposta correta, confiando na instrução de autoconferência dada ao
modelo. Uma auditoria manual de amostra dos itens gerados por nível, e
eventualmente uma verificação simbólica automatizada para os tópicos em que ela
é viável, são consideradas pré-requisito para a coleta em campo. A segunda é que
o limitador de taxa mantém estado em memória do processo, o que basta para uma
instância única, mas exigiria um contador compartilhado em uma implantação
distribuída. A terceira é a ausência de cobertura de testes automatizados: a
validação foi manual, e a inclusão de testes sobre o motor de adaptabilidade,
sobre a validação da questão gerada e sobre o filtro anti-vazamento é
considerada requisito antes de qualquer coleta de dados com estudantes reais.

% ================================================================
\section{Conclusão Esperada}
\label{sec:conclusao}

A hipótese central deste trabalho é que o auxílio mediado socrático produz maior
retenção de conhecimento do que o fornecimento direto de respostas, ainda que ao
custo de maior esforço e maior tempo de resolução durante a intervenção. Espera-se,
portanto, o seguinte conjunto de resultados.

\textbf{H1 --- Ganho de aprendizagem.} Espera-se que o ganho normatizado médio do
grupo socrático supere o do grupo diretivo, com diferença estatisticamente
significativa a $\alpha = 0{,}05$. A magnitude esperada é modesta: dada a duração
limitada da intervenção e o tamanho amostral, um tamanho de efeito na faixa
pequena a moderada já seria consistente com a literatura de
\textit{scaffolding}~\cite{kestin2025,ocen2026}.

\textbf{H2 --- Desempenho durante a intervenção.} Espera-se que o grupo diretivo
apresente acurácia igual ou superior durante a fase de exercícios, e menor número
de tentativas até o acerto. Esse é o resultado esperado e não contradiz H1: é
precisamente a manifestação do descarregamento cognitivo, em que o desempenho
imediato melhora sem correspondência na retenção posterior. A divergência entre
desempenho durante a intervenção e desempenho no pós-teste, se confirmada, é o
achado mais informativo do estudo.

\textbf{H3 --- Viabilidade operacional.} Espera-se que a latência mediana das
chamadas ao modelo permaneça em faixa compatível com uso em sala de aula e que a
taxa de vazamento detectado no modo socrático se mantenha baixa. Uma taxa de
vazamento elevada não invalidaria apenas o produto, mas o próprio experimento, e
por isso é tratada como métrica de validade interna.

Espera-se ainda que a contribuição de engenharia se sustente
independentemente do resultado estatístico: o conjunto de salvaguardas descrito
--- retenção do gabarito no servidor, autorrelato estruturado, filtro heurístico,
validação estrutural dos itens gerados e \textit{fallback} pedagógico em ambos os
caminhos --- constitui um padrão reaplicável a qualquer
estudo que pretenda contrastar modos de mediação de IA em contexto educacional,
onde a separação entre os tratamentos não pode ser presumida a partir do
\textit{prompt}.

Quanto aos resultados que refutariam a hipótese, três cenários são
considerados. Ausência de diferença significativa entre os grupos indicaria que,
nesta duração e neste domínio, a forma de mediação não é determinante --- o que
seria um achado relevante contra o entusiasmo corrente com \textit{prompts}
socráticos. Ganho superior do grupo diretivo sugeriria que, para alunos em
defasagem acentuada, a exposição direta ao raciocínio correto é mais eficaz que a
mediação, hipótese compatível com a noção de sobrecarga cognitiva em estudantes
sem repertório prévio suficiente. E uma taxa alta de vazamento tornaria o
contraste inconclusivo, exigindo reanálise restrita às interações íntegras.

As limitações antecipadas são a amostra reduzida e de conveniência, a restrição
a um único domínio de conteúdo, a curta duração da intervenção --- que não
permite avaliar retenção de longo prazo --- e a impossibilidade de cegamento dos
participantes quanto ao tipo de feedback recebido, ainda que desconheçam a
existência do outro grupo. Soma-se a elas uma limitação introduzida pela própria
geração automática dos itens: como cada aluno responde a um conjunto de questões
inédito e distinto, os grupos não são comparados sobre exatamente os mesmos
estímulos durante a intervenção. O contraste principal permanece válido, porque o
pré-teste e o pós-teste são fixos e idênticos para todos, e porque a geração é
cega ao grupo experimental; ainda assim, a acurácia medida durante a fase de
exercícios passa a carregar variância adicional atribuível à dificuldade
heterogênea dos itens gerados, e deve ser lida com essa ressalva. Como
desdobramentos, projetam-se a verificação automatizada do gabarito dos itens
gerados, a inclusão do módulo de Língua Portuguesa, a adição de pós-teste tardio
para medir retenção e a evolução do painel docente para acompanhamento individual
em tempo real.

% ================================================================
\section{Cronograma Geral de Execução}
\label{sec:cronograma}

O ciclo de atividades foi distribuído ao longo do primeiro semestre de 2026,
conforme a Tabela~\ref{tab:cronograma}. O levantamento bibliográfico e a
engenharia de requisitos ocorrem em fevereiro e março, em paralelo à auditoria
dos itens gerados pela IA e à inclusão de testes automatizados sobre o motor de
adaptabilidade e o filtro anti-vazamento. A coleta em campo e o tratamento
estatístico dos dados concentram-se em abril e maio, junto à escrita final. A
defesa perante a banca examinadora está prevista para junho.

\begin{table}[H]
\centering
\caption{Cronograma de execução do projeto (2026)}
\label{tab:cronograma}
\vspace{0.13cm}
\resizebox{\columnwidth}{!}{
\begin{tabular}{lccccc}
\toprule
\textbf{Atividade} & \textbf{Fev} & \textbf{Mar} & \textbf{Abr} & \textbf{Mai} & \textbf{Jun} \\ \midrule
1. Levantamento e requisitos      & X & X &   &   &   \\
2. Auditoria dos itens gerados    & X & X &   &   &   \\
3. Testes automatizados           &   & X &   &   &   \\
4. Coleta em campo (TCLE/TALE)    &   &   & X &   &   \\
5. Tratamento estatístico         &   &   & X & X &   \\
6. Escrita final                  &   &   &   & X &   \\
7. Defesa perante a banca         &   &   &   &   & X \\ \bottomrule
\end{tabular}%
}
\end{table}

% ================================================================
\begin{thebibliography}{9}
\small

\bibitem{ocen2026}
OCEN, P. et al.
\textbf{Artificial intelligence as a \textit{scaffolding} tool for self-directed learning}.
\textit{Frontiers in Education}, 2026.
Disponível em: \url{https://www.frontiersin.org/journals/education/articles/10.3389/feduc.2026.1793940/full}.
Acesso em: maio 2026.

\bibitem{kestin2025}
KESTIN, G. et al.
\textbf{Scaffolding Generative AI as a Tutor: A Quasi-Experimental Study}.
\textit{MDPI Education}, 2025.
Disponível em: \url{https://www.mdpi.com/2227-7102/16/4/651}.
Acesso em: maio 2026.

\bibitem{roe2025}
VOROBYEVA, Klarisa I. et al.
\textbf{Personalized learning through AI: Pedagogical approaches and critical insights}.
\textit{Contemporary Educational Technology}, 2025.
Disponível em: \url{https://www.cedtech.net/article/personalized-learning-through-ai-pedagogical-approaches-and-critical-insights-16108}.
Acesso em: maio 2026.

\bibitem{kim2025}
KIM, J. et al.
\textbf{Students--generative AI interaction patterns and its impact on academic writing}.
\textit{Journal of Computing in Higher Education}, 2025.
DOI: \url{https://doi.org/10.1007/s12528-025-09444-6}.

\bibitem{pmc2026}
PMC.
\textbf{Evaluating AI-powered learning assistants in engineering higher education}.
\textit{PubMed Central}, 2026.
Disponível em: \url{https://pmc.ncbi.nlm.nih.gov/articles/PMC12932848/}.
Acesso em: maio 2026.

\bibitem{cetic2025}
BARBOSA, Alexandre F.
\textbf{Estudos Setoriais -- Inteligência Artificial na Educação}. 2025.
Disponível em: \url{https://cetic.br/media/docs/publicacoes/7/pt/20251124074142/estudos_setoriais-ia_na_educacao.pdf}.
Acesso em: maio 2026.

\bibitem{mec2023}
BRASIL. Ministério da Educação (MEC).
\textbf{Estratégia Nacional de Escolas Conectadas (Enec)}. 2023.
Disponível em: \url{https://www.gov.br/mec/pt-br/escolas-conectadas}.
Acesso em: maio 2026.

\bibitem{luckin2018}
LUCKIN, R.
\textbf{Machine Learning and Human Intelligence}.
London: UCL Press, 2018.
Disponível em: \url{https://discovery.ucl.ac.uk/id/eprint/10051189/}.
Acesso em: maio 2026.

\bibitem{unesco2024}
UNESCO.
\textbf{Guidance for Generative AI in Education and Research}.
Paris: UNESCO Publishing, 2024.

\end{thebibliography}

\end{document}
